\section{Two ramified mesh families with vanishing classes}
\label{sec:pure-power-vanishing}

For the two coefficient rings considered here, the standard mesh
with determinant \(g(X)+t^n\) has the regular-lattice class at every
vertex.  We prove this using explicit coefficient-split sequences.
Polynomial generation nevertheless fails for these rings by
Theorem~\ref{thm:ramified-vertical-truncations}.  The calculation
therefore identifies a specific family of zero classes within those
negative cases.

Throughout the mesh argument, let \(O\) be a complete
characteristic-zero DVR with residue characteristic \(p\), and put
\[
 A=O[t]/(t^n),\qquad P=C_p=\langle\sigma\rangle,
 \qquad R=A[P],\qquad X=\sigma-1,
\]
where \(n\ge2\).  Write
\[
 g(X)=(1+X)^p-1=Xh(X),\qquad
 h(X)=p+\binom p2X+\cdots+X^{p-1}.
\]
For \(0\le i\le n\), let \(M_i\) have generators \(a,b\) and
presentation
\begin{equation}\label{eq:pure-power-mesh-presentation}
 \Phi_i=\begin{pmatrix}X&t^i\\-t^{n-i}&h(X)\end{pmatrix},
 \qquad Xa=t^{n-i}b,\qquad t^ia+h(X)b=0.
\end{equation}

\begin{theorem}\label{thm:pure-power-mesh-vanishing}
In each of the following cases,
\[
 \begin{array}{c|c|c}
 p&n&O\\ \hline
 3&3&\mathbb Z_3[\pi],\quad\pi^3=3\\
 2&4&\mathbb Z_2[\pi],\quad\pi^4=2,
 \end{array}
\]
one has
\[
 [M_i]=[A[C_p]]\quad\text{in }K_0(\Lat_A(C_p))
 \qquad(0\le i\le n).
\]
Their induced differences also vanish in
\(K_0(\Lat_A(S_p))\).
\end{theorem}

The displayed rings are the complete DVRs defined by the indicated
Eisenstein polynomials.  The theorem concerns these particular
presentations and their induction images.

\subsection{The coefficient-split mesh}

The determinant of \eqref{eq:pure-power-mesh-presentation} is
\(g+t^n\), so its cokernel over \(A\) is annihilated by \(g\).
It is \(A\)-free of rank \(p\), with basis
\[
 a,\qquad X^jb\quad(0\le j\le p-2).
\]
The first relation specifies \(Xa\), and the monic second relation
specifies \(X^{p-1}b\).  Conversely, these prescriptions define
the operator on the displayed free module and satisfy both relations,
giving the asserted presentation.  At \(i=0\) eliminate \(a=-hb\);
at \(i=n\) eliminate \(b=Xa\).  The remaining relation is \(g=0\),
so \(M_0\simeq M_n\simeq R\).

In these bases, the maps
\[
 \alpha_i:M_i\longrightarrow M_{i+1},\quad
 \alpha_i=\operatorname{diag}(t,I_{p-1}),\qquad
 \beta_i:M_{i+1}\longrightarrow M_i,\quad
 \beta_i=\operatorname{diag}(1,tI_{p-1})
\]
are equivariant by substitution in the relations.  Their composites
are multiplication by \(t\).  For \(1\le i<n\), they give
\begin{equation}\label{eq:pure-power-mesh}
 0\longrightarrow M_i
 \xrightarrow{(\beta_{i-1},\alpha_i)}M_{i-1}\oplus M_{i+1}
 \xrightarrow{(\alpha_{i-1},-\beta_i)}M_i\longrightarrow0.
\end{equation}
On coefficient vectors \((r,v)\), the maps are
\[
 (r,v)\longmapsto((r,tv),(tr,v)),\qquad
 ((r_L,v_L),(r_R,v_R))\longmapsto(tr_L-r_R,v_L-tv_R).
\]
The kernel equations give exactness.  A coefficient left inverse
of the first map and a coefficient right inverse of the second are
\[
 ((r_L,v_L),(r_R,v_R))\longmapsto(r_L,v_R),\qquad
 (r,v)\longmapsto((0,v),(-r,0)).
\]
Thus \eqref{eq:pure-power-mesh} is admissible.  Setting
\(x_i=[M_i]-[R]\), the mesh and endpoints give
\begin{equation}\label{eq:pure-power-mesh-relations}
 x_i=i x_1,\qquad nx_1=0,\qquad x_{n-1}=-x_1.
\end{equation}

\subsection{An elementary finite-module criterion}

Set
\[
 C=O[[t]],\qquad S=C[X]/(g(X)+t^n),\qquad J=(X,t)\subset S.
\]
The ring \(S\) is \(C\)-free with basis \(1,X,\ldots,X^{p-1}\),
and \(S/t^nS=R\).  Moreover \(J\) is \(C\)-free with basis
\(t,X,\ldots,X^{p-1}\): reducing \(X^p\) produces a constant
multiple of \(t^n\), and these displayed elements are independent
in the \(C\)-basis of \(S\).
The relations on this basis give
\begin{equation}\label{eq:pure-power-ideal-lift}
 J\simeq\operatorname{coker}
 \begin{pmatrix}X&t^{n-1}\\-t&h(X)\end{pmatrix},
 \qquad a\longmapsto t,\quad b\longmapsto X,
 \qquad J/t^nJ\simeq M_{n-1}.
\end{equation}
Write \(\varepsilon:S\to S/J=O\) for evaluation at \(t=X=0\).

\begin{proposition}\label{prop:pure-power-finite-module-criterion}
Suppose a finite \(S\)-module \(N\), free over \(C\), has an
\(S\)-linear endomorphism \(b\) fitting into an exact sequence
\begin{equation}\label{eq:pure-power-normalization-sequence}
 0\longrightarrow N\xrightarrow{b}N
 \xrightarrow{\nu}O\longrightarrow0,
\end{equation}
where \(S\) acts on \(O\) through \(\varepsilon\).
Then \([M_i]=[R]\) in \(K_0(\Lat_A(C_p))\) for every
\(0\le i\le n\).
\end{proposition}

\begin{proof}
Form the finite \(S\)-module
\[
 E=S\times_O N
   =\{(s,u):\varepsilon(s)=\nu(u)\}.
\]
Its two projections give exact sequences
\begin{equation}\label{eq:pure-power-fiber-product-sequences}
 \begin{split}
 0&\longrightarrow J\longrightarrow E\longrightarrow N
       \longrightarrow0,\qquad j\longmapsto(j,0),\\
 0&\longrightarrow N\longrightarrow E\longrightarrow S
       \longrightarrow0,\qquad u\longmapsto(0,b(u)).
 \end{split}
\end{equation}
The second kernel identification uses
\(\ker\nu=\operatorname{im}b\) and the injectivity of \(b\).

Here are coefficient splittings.  Choose a \(C\)-basis
\(u_1,\ldots,u_d\) of \(N\), put \(a_j=\nu(u_j)\in O\),
and define
\[
 \lambda\!\left(\sum_j f_j(t)u_j\right)
       =\sum_j f_j(t)a_j\in C\subset S.
\]
Then \(\varepsilon\lambda=\nu\), so
\(u\mapsto(\lambda(u),u)\) splits \(E\to N\) over \(C\).
Choose \(u_*\in N\) with \(\nu(u_*)=1\).  If
\(s=\sum_{j=0}^{p-1}s_j(t)X^j\), then
\(s\mapsto(s,s_0(t)u_*)\) is a \(C\)-linear section of
\(E\to S\).  Thus both sequences are coefficient-split and
\(E\) is \(C\)-free.

Reduction modulo \(t^n\) preserves these splittings.  Using
\eqref{eq:pure-power-ideal-lift}, write
\(\overline E=E/t^nE\) and \(\overline N=N/t^nN\).
The resulting admissible sequences in \(\Lat_A(C_p)\) are
\[
 \begin{gathered}
 0\longrightarrow M_{n-1}\longrightarrow\overline E
       \longrightarrow\overline N\longrightarrow0,\\
 0\longrightarrow\overline N\longrightarrow\overline E
       \longrightarrow R\longrightarrow0.
 \end{gathered}
\]
Their relations give \([M_{n-1}]=[R]\).  Now
\eqref{eq:pure-power-mesh-relations} gives \(x_1=0\), and then
\(x_i=0\) for every \(i\).
\end{proof}

\subsection{The two finite algebras}

For \(p=n=3\) and \(\pi^3=3\), define
\begin{equation}\label{eq:pure-power-cubic-polynomial}
 P_3(v,t)=v^3+\pi^2(1-t)v^2+\pi(1-t)^2v+t(t-1).
\end{equation}
Direct expansion gives
\begin{equation}\label{eq:pure-power-cubic-identity}
 g(\pi v-t)+t^3=\pi^3P_3(v,t).
\end{equation}
Consequently \(N=C[v]/(P_3)\), which is \(C\)-free of rank three,
is an \(S\)-algebra by \(X\mapsto\pi v-t\).
Multiplication by \(v\) has determinant, up to sign,
\(P_3(0,t)=t(t-1)\ne0\).  Since \(C\) is a domain and \(N\)
is \(C\)-free, this multiplication is injective.  Also
\[
 N/vN=C/(t(t-1))=C/(t)=O,
\]
with \(X=t=0\) on the quotient.  Thus
\begin{equation}\label{eq:pure-power-cubic-normalization}
 0\longrightarrow N\xrightarrow{v}N
       \longrightarrow S/J\longrightarrow0
\end{equation}
satisfies the finite-module criterion.

For \(p=2,n=4\) and \(\pi^4=2\), put
\[
 q(t)=-\pi^2t-(1+\pi^2)t^2
\]
and define
\begin{equation}\label{eq:pure-power-quartic-polynomial}
 P_4(z,t)=z^2+\pi(1+q(t))z
       +t(1+t)\bigl((1+\pi^2)t^2-1\bigr).
\end{equation}
The substitution
\begin{equation}\label{eq:pure-power-quartic-coordinate}
 X=\pi^3z+q(t)
\end{equation}
gives the identity
\begin{equation}\label{eq:pure-power-quartic-identity}
 g(\pi^3z+q(t))+t^4=\pi^6P_4(z,t).
\end{equation}
Indeed its linear coefficient in \(z\) is
\(2\pi^3(1+q)=\pi^7(1+q)\), and its constant coefficient is
\[
 q(t)^2+2q(t)+t^4
       =\pi^6t(1+t)\bigl((1+\pi^2)t^2-1\bigr).
\]
Thus \(N=C[z]/(P_4)\) is a \(C\)-free \(S\)-algebra of rank two.
Its constant term
\begin{equation}\label{eq:pure-power-quartic-constant}
 P_4(0,t)=t(1+t)\bigl((1+\pi^2)t^2-1\bigr)
\end{equation}
is \(t\) times a unit of \(C\).  The multiplication determinant
therefore proves that \(z\) is injective on \(N\), and
\(N/zN\simeq O\), with \(X=t=0\).  We obtain
\begin{equation}\label{eq:pure-power-quartic-normalization}
 0\longrightarrow N\xrightarrow{z}N
       \longrightarrow S/J\longrightarrow0.
\end{equation}
The criterion applies again and proves
Theorem~\ref{thm:pure-power-mesh-vanishing}.  Induction to \(S_p\)
preserves coefficient-free modules and coefficient-split sequences,
and therefore preserves these zero differences.

\subsection[A surface divisor for the unit one plus pi]{A surface divisor for the unit \(1+\pi\)}

The following proposition describes a different integral parameter.
It is a statement about a normal surface and its divisor class group;
the exact-category mesh theorem above has already been proved.

\begin{proposition}\label{prop:pure-power-unit-surface}
Let \(n=p^a\), \(a\ge1\), and let \(O\) be a complete
characteristic-zero DVR with uniformizer \(\pi\) and
\(v_\pi(p)\ge2\).  Set \(c=1+\pi\) and
\[
 S_c=O[[X,t]]/(g(X)+ct^n),\qquad J_c=(X,t).
\]
Then \(S_c\) is a normal two-dimensional local domain, and
\([J_c]\) has exact order \(n\) in \(\operatorname{Cl}(S_c)\).
\end{proposition}

\begin{proof}
As a monic polynomial in \(t\) over \(O[[X]]\), the equation
\(t^n+Xh(X)/c\) is Eisenstein at \((X)\): the element \(h(0)=p\)
is a unit in \(O[[X]]_{(X)}\).  The quotient is finite free over
\(O[[X]]\), so irreducibility over its fraction field proves that
\(S_c\) is a domain.  It is a hypersurface of dimension two and
is Cohen--Macaulay.

Put \(r=p^{a-1}\) and \(Y=X+t^r\).  Modulo \(\pi\), the defining
equation is \(Y^p\).  Thus the unique height-one prime over
\(\pi\) is the image of \((\pi,Y)\).  In the regular ambient
ring localized at this prime, \(t\) is a unit, and
\begin{equation}\label{eq:pure-power-unit-vertical}
 g(X)+(1+\pi)t^n\equiv\pi t^n\pmod{(\pi,Y)^2}.
\end{equation}
All intermediate binomial coefficients lie in \((p)\subseteq
(\pi^2)\).  The pure \(t^n\) terms cancel for odd \(p\); at
\(p=2\), their sum is \(2t^n\in(\pi^2)\).
The right side of \eqref{eq:pure-power-unit-vertical} is outside
the square of the ambient maximal ideal, so this height-one
localization is regular.

At a height-one prime not containing \(\pi\), the derivative
\(g'(X)=p(1+X)^{p-1}\) is a unit.  Here \(1+X\) is already a
unit in the formal ambient ring.  These localizations are therefore
regular as well.  Serre's \(R_1+S_2\) criterion proves normality
\citep[Tag 031S]{Stacks2026}.

For \(1\le b\le n\), put \(I_b=(X,t^b)\).
Its quotient \(S_c/I_b\simeq O[t]/t^b\) has depth one, so the
depth lemma gives depth two for \(I_b\).  At height-one primes
the ideal is torsion-free over a DVR.  Hence it is a rank-one
reflexive module, and represents a divisor class
\citep[Tags 0AVB and 0EBM]{Stacks2026}.

The prime \(J_c\) has quotient \(O\).  At its localization,
\(p\), \(h(X)\), and \(c\) are units, so \(Xh=-ct^n\) gives
\[
 v_{J_c}(t)=1,\qquad v_{J_c}(X)=n.
\]
It is the only height-one prime containing \(X\).  Consequently
\[
 \operatorname{div}(X)=nJ_c,\qquad
 [I_b]=b[J_c]\quad\text{in }\operatorname{Cl}(S_c),
\]
using a consistent ideal-class convention.

For \(1\le b<n\), the element \(t^b\) is nonzero modulo \(X\),
since that quotient is \(O[t]/t^n\).  Also \(X\notin(t^b)\):
setting \(t=0\) leaves the nonzero class of \(X\) in
\(O[[X]]/(Xh)\), because \(h\) is a nonunit.
If \((X,t^b)\) were principal, write \(X=ur\), \(t^b=us\),
and \(u=\alpha X+\beta t^b\) for a generator \(u\ne0\).
Then \(1=\alpha r+\beta s\), so one of \(r,s\) is a unit in
the local ring.  One of \(X,t^b\) would generate the ideal,
contrary to the two preceding observations.  Thus every \(I_b\),
\(1\le b<n\), is nonprincipal.  Together with
\(\operatorname{div}(X)=nJ_c\), this proves exact order \(n\).
\end{proof}

\begin{remark}[Scope of the surface calculation]
\label{rem:pure-power-missing-specialization}
The failure of polynomial generation for the two rings in
Theorem~\ref{thm:pure-power-mesh-vanishing} follows from
Theorem~\ref{thm:ramified-vertical-truncations}.  The separate
question addressed here is whether the fixed surface class of
Proposition~\ref{prop:pure-power-unit-surface} itself yields a
whole-lattice detector through the corresponding parameter chart.
The proposition alone does not construct that specialization functor.

The completed family at this parameter is
\[
 O[[X,t,w]]/(g(X)+(c+w)t^n).
\]
Trivializing it by an integral rescaling of \(t\) would require
\(u(w)^n=1+w/c\) with \(u(w)\in1+wO[[w]]\).  Its linear
coefficient would be \(1/(nc)\), which is not integral.
Derived restriction of a module killed by \(w\) is the
two-term zero-differential complex on that module and has zero
Grothendieck class.

The ramified substitution \(w=c((1+v)^n-1)\) does trivialize the
family after a finite free cover of degree \(n\).  The inverse
image of \(w=0\) has coefficient ring
\[
 O[[v]]/((1+v)^n-1),
\]
free of rank \(n\) over \(O\).  Replacing \(t\) by the unit
multiple \((1+v)t\) preserves the ideal \((X,t^b)\); forgetting
\(v\) therefore gives \(n\) copies of that ideal.
In the Grothendieck group of
\(D_{\mathrm{sg}}(S_c)=D^b(\operatorname{mod}S_c)/
\operatorname{Perf}(S_c)\), the \(t\)-filtration of
\(S_c/(X,t^b)\) gives \([I_b]=b[J_c]\).
Since \(I_n=(X)\simeq S_c\), one has \(n[J_c]=0\).
Thus this cover kills the displayed torsion witness.  These
calculations delimit the two proposed specialization procedures;
they do not change the negative generation conclusion cited above.
\end{remark}
